\section{Multiple Qubits}
Quando si vanno a considerare sistemi con più qubit, allora entra in gioco il \textbf{prodotto tensoriale}.
Tale operazione si distingue dalla somma, prima di passare alla sua formalizzazione matematica è importante fare una considerazione sulla dimensione dello spazio vettoriale complesso in cui si andrà a lavorare. Con la somma tra due spazi vettoriali sì ottiene la seguente condizione:
\[
\mathbb{C}^n \oplus \mathbb{C}^m = \mathbb{C}^{n+m}
\]
Dove la base viene definira mediante l'unione delle due basi del sottospazio:
\[
\mathcal{B} = \{v_1,\dots,v_n, u_1, \dots, u_m\}
\]
Mentre con il prodotto tensoriale, lo spazio sarà di dimensione:
\[
\mathbb{C}^n \otimes \mathbb{C}^m = \mathbb{C}^{nm}
\]
Dove la base sarà costituita da:
\[
\mathcal{B} = \{v_1 \otimes u_1, v_1 \otimes u_2, \dots v_n \otimes u_m\}
\]

Se andiamo a considerare in maniera diretta due stazi: \(V\) e \(W\), associati a due qubit, questi saranno dotati di due basi ben distinte:
\[
\mathcal{B}_V = \{v_1, v_2\}
\]
\[
\mathcal{B}_W = \{w_1, w_2\}
\]

Dove, la base intesa per la somma degli spazi, sarà:
\[
\mathcal{B}_{V \oplus W}= \{v_1,v_2,w_1,w_2\}
\]

Mentre per il prodotto vettoriale, si avrà:
\[
\mathcal{B}_{V \otimes W} = \{v_1 \otimes w_1, v_1 \otimes w_2, v_2 \otimes w_1, v_2 \otimes w_2 \}
\]

Quando si va a considerare la base per il sistema di interazione tradue qubit, espressi nella base canonica, allora si avrà:
\[
\mathcal{B} = \{\ket{0} \otimes \ket{0}, \ket{0} \otimes \ket{1}, \ket{1} \otimes \ket{0}, \ket{1} \otimes \ket{1}\}
\]

Che più comunenemente, tale base, viene scritta come:
\[
\mathcal{B} = \{\ket{00}, \ket{01}, \ket{10}, \ket{11}\}
\]

Andando a fare un confronto con quello che è il caso classico, il numero di combinazioni con i qubit è uguale al numero di combinazioni possibili all'interno dei sistemi classici.

\begin{info}
Le posizioni all'interno del vettore con la ket, non perdono il loro significato funzionale. Quindi:
\[
\ket{00} = \ket{0_A 0_B}
\]
Dove \(A\) e \(B\), sono due qubit diversi e ben identificati. Il non specificare a quale sistema stiamo facendo riferimento è un abuso di notazione ( opiù una mancanza di notazione)
\end{info}

La differenza sostanziale sta nel fatto che nel caso classico, in un determinato istante, un registro può trovarsi in una sola delle sue combinazioni, mentre nel caso quantistico, una coppia di qubit ha una probabilità non nulla di potersi trovare in uno dei 4 possibili stati della base (dopo la lettura).


Detto questo, quindi, è possibile scrivere un qubit appartenente ad un sistema di 2 qubit come:
\[
\ket{\psi} 0 \alpha_0 \ket{00} + \alpha_1 \ket{01} + \alpha_2 \ket{10} + \alpha_3 \ket{11}
\]
\[
\alpha_0, \alpha_1, \alpha_2, \alpha_3 \in \mathbb{C}, |\alpha_0|^2 + |\alpha_1|^2 + |\alpha_2|^2 + |\alpha_3|^2 = 1
\]

Andando poi a generalizzare la cosa, per un sistema ad n-qubit, si avrà che:
\[
\ket{psi} = \sum_{i=0}^{2^n-1} \alpha_i \ket{i}
\]
\[
\alpha_i \in \mathbb{C} : \sum_{i}|\alpha_i|^2=1
\]